\chapter{Locality, Continuous States, and Momentum}

\lectureoverview{Entanglement can be created and undone by reversible dynamics, yet measurements become effectively irreversible when part of the entangled system escapes our control. That distinction leads to a precise locality test: Bob may apply any unitary operation in his neighborhood, but Alice's reduced density matrix, and therefore every statistic available to her alone, remains unchanged. A classical-simulation thought experiment then isolates what is genuinely nonclassical. One spin is easy to imitate, and a central computer can imitate an entangled pair, but separated local machines cannot reproduce all entangled correlations without illicit communication or a Bell-local hidden-variable model that ultimately fails. The lecture then leaves spins for a particle on a line. Discrete sums become integrals, Kronecker deltas become Dirac deltas, and wave functions become vectors in an infinite-dimensional space. Multiplication by \(x\) gives position; differentiation is anti-Hermitian; multiplying it by \(-i\) produces the Hermitian momentum operator. The final calculation, \([x,p]=i\hbar\), reconnects quantum commutators with classical Poisson brackets and points toward wave packets, Fourier transforms, and Schr\"odinger evolution.}

\section{Entanglement Is Not Irreversibility}

Begin with a distinction that is easy to lose in discussions of measurement. Suppose two spins arrive in a product state, interact, separate, and leave in an entangled state. Must the entanglement remain forever? No. The interaction was unitary, so another suitable interaction can undo it. One may reflect the spins, bring them together again, and arrange an inverse evolution that returns the pair to a product state. There is no law according to which entanglement itself only increases.

The situation changes operationally when one member of the pair is discarded. If an electron flies away into an uncontrolled environment, then the partner that remains nearby is still entangled with something, but the missing degrees of freedom are no longer available for the inverse operation. In principle the distant electron might be reflected back; in practice it is gone. This is the sense in which a measurement is irreversible \emph{for us}: the correlations have spread into degrees of freedom that we do not control.

\begin{classroomqa}
\audiencequestion{If two systems become entangled, can they later become disentangled?}
\lecturerresponse{Yes. Unitary evolution is reversible. A second interaction can implement the inverse of the first and restore a product state. Irreversibility enters when relevant degrees of freedom are thrown away or absorbed into an uncontrolled environment, not from entanglement alone.}
\end{classroomqa}

\begin{classroomqa}
\audiencequestion{Then why does a completed measurement seem irreversible?}
\lecturerresponse{Because the measured system has become correlated with an apparatus, an observer, and ultimately a large environment. Reversing the complete state would require control of all those records. The microscopic evolution may remain unitary while the operation is inaccessible from the observer's reduced description.}
\end{classroomqa}

\section{The State Available to Alice}

Let a composite system be divided into Alice's degrees of freedom \(a\) and Bob's degrees of freedom \(b\), with pure-state coefficients
\begin{equation}
\Psi(a,b)=\braket{a,b}{\Psi}.
\end{equation}
Alice does not have access to Bob's laboratory. Her complete local description is therefore not generally a wave function. It is the reduced density matrix
\begin{equation}
\rho_A(a',a)
=
\sum_b \Psi^*(a',b)\Psi(a,b).
\label{eq:l8-reduced-density}
\end{equation}
The sum removes Bob's label. What remains has one bra index and one ket index for Alice, and it contains every probability distribution that Alice can obtain by experiments performed entirely in her neighborhood.

\begin{figure}[htbp]
\centering
\includegraphics[width=0.88\textwidth]{lecture_08_figure_01.png}
\caption{Alice's reduced density matrix, formed by summing the joint wave function over Bob's basis label.}
\lectureframe{00:04:34}
\end{figure}

Here Alice's neighborhood and Bob's neighborhood should be read broadly. Each may include apparatus, magnetic fields, technicians, and any other degrees of freedom with which that observer can interact. The claim of locality is not that Bob cannot change his own system. He may do so in many ways. The claim is that, while he remains distant, nothing he does locally changes Equation~\eqref{eq:l8-reduced-density}.

\section{Bob's Local Unitary}

First ignore Alice and describe Bob's manipulation alone. If Bob changes a magnetic field or otherwise changes his local Hamiltonian, his wave-function coefficients undergo a unitary transformation,
\begin{equation}
\psi_M(b)
=
\sum_{b'}U_{bb'}\psi(b').
\label{eq:l8-bob-unitary}
\end{equation}
In the composite language Alice's index is a spectator:
\begin{equation}
\Psi_M(a,b)
=
\sum_{b'}U_{bb'}\Psi(a,b').
\label{eq:l8-composite-unitary}
\end{equation}
Complex conjugation converts the ket transformation into the corresponding bra transformation,
\begin{equation}
\Psi_M^*(a',b)
=
\sum_{b''}\Psi^*(a',b'')U^\dagger_{b''b}.
\label{eq:l8-composite-bra}
\end{equation}

\begin{figure}[htbp]
\centering
\includegraphics[width=0.88\textwidth]{lecture_08_figure_02.png}
\caption{A unitary acts on Bob's ket index, while the Hermitian-conjugate matrix acts on the corresponding bra index.}
\lectureframe{00:12:54}
\end{figure}

Now recalculate Alice's density matrix:
\begin{align}
\rho_{A,M}(a',a)
&=
\sum_b\Psi_M^*(a',b)\Psi_M(a,b) \notag\\
&=
\sum_{b,b',b''}
\Psi^*(a',b'')
U^\dagger_{b''b}U_{bb'}
\Psi(a,b') \notag\\
&=
\sum_{b',b''}
\Psi^*(a',b'')
(U^\dagger U)_{b''b'}
\Psi(a,b') \notag\\
&=
\sum_{b',b''}
\Psi^*(a',b'')\delta_{b''b'}\Psi(a,b') \notag\\
&=
\sum_{b'}\Psi^*(a',b')\Psi(a,b') \notag\\
&=
\rho_A(a',a).
\label{eq:l8-locality-proof}
\end{align}

\begin{figure}[htbp]
\centering
\includegraphics[width=0.88\textwidth]{lecture_08_figure_03.png}
\caption{The locality calculation on the board: \(U^\dagger U\) collapses to the identity inside Alice's reduced state.}
\lectureframe{00:14:27}
\end{figure}

The entire proof is the contraction
\begin{equation}
\sum_bU^\dagger_{b''b}U_{bb'}
=(U^\dagger U)_{b''b'}
=\delta_{b''b'}.
\label{eq:l8-unitarity-contraction}
\end{equation}
No factorization assumption was used. Equation~\eqref{eq:l8-locality-proof} holds whether the joint state is a product or is maximally entangled. Thus every Alice-only expectation value,
\begin{equation}
\expect{L_A}
=
\operatorname{Tr}_A(\rho_A L_A),
\end{equation}
is invariant under Bob's local unitary.

\begin{classroomqa}
\audiencequestion{If Bob is doing something to his spin, why is that not already a violation of unitarity?}
\lecturerresponse{Doing something does not mean applying a nonunitary map. Bob can rotate a magnet, change a local Hamiltonian, or let his spin precess. Those are ordinary unitary evolutions. They alter his own state and his own outcome probabilities, but the calculation above shows that they do not alter Alice's reduced state.}
\end{classroomqa}

\begin{classroomqa}
\audiencequestion{Does quantum mechanics therefore obey locality?}
\lecturerresponse{It obeys the operational locality used here: Bob cannot influence the statistics of measurements available to Alice by acting only in his neighborhood. Entangled joint correlations remain nonclassical, but they do not provide a controllable signal.}
\end{classroomqa}

\subsection{Why unitarity matters}

Unitarity was not decorative. It supplied the identity in Equation~\eqref{eq:l8-unitarity-contraction}. In earlier language, unitarity expresses the conservation of distinctions: two distinguishable state vectors remain distinguishable under reversible evolution. If Bob possessed even a slightly nonunitary deterministic evolution, \(U^\dagger U\) would not be the identity. Alice's reduced density matrix could then change immediately, giving Bob a means of signaling faster than light.

This tight connection is one reason quantum mechanics is difficult to modify piecemeal. Change the evolution law and apparently remote principles can fail with it. Locality, probability conservation, and reversible state evolution are parts of one structure.

\begin{classroomqa}
\audiencequestion{How would a nonunitary operation permit a faster-than-light signal?}
\lecturerresponse{Bob would choose whether or not to apply it. If the map changed \(\rho_A\), Alice could distinguish the two choices by repeating local measurements on an ensemble. The change would occur without a signal propagating from Bob, which is precisely what the unitary cancellation prevents.}
\end{classroomqa}

\section{Collapse Changes Conditions, Not Distant Marginals}

The beam-splitter example makes the uneasiness vivid. A single photon is placed in a superposition of two paths, one toward Alice and one toward Bob. If Alice's detector clicks, the conditional state assigns zero amplitude to a later click on Bob's path. It is tempting to say that Alice has instantaneously changed Bob's photon.

What Bob can observe without Alice's record is different. Before receiving any message, he must average over Alice's possible outcomes. That unconditioned average is described by his reduced density matrix, and it does not change. Alice's click changes which conditional subensemble describes a run once her result is known; it does not change Bob's local marginal distribution. The distinction between
\begin{equation}
\rho_B
\qquad\text{and}\qquad
\rho_{B\mid A=\text{click}}
\end{equation}
is the distinction between no signaling and useful correlations.

\begin{classroomqa}
\audiencequestion{When one detector clicks, does the wave function on the other arm not collapse to zero instantaneously?}
\lecturerresponse{The state conditioned on that click changes, but the distant observer does not possess the condition until an ordinary message arrives. Bob's unconditioned density matrix and all of his average local statistics remain unchanged. Collapse cannot be used as a telegraph.}
\end{classroomqa}

\begin{classroomqa}
\audiencequestion{Is the confusion partly a confusion between the state and information about the state?}
\lecturerresponse{That wording is suggestive but still too loose. The concrete distinction is between a joint state, a reduced state, and a state conditioned on a reported outcome. Keeping those three mathematical objects separate is safer than trying to settle the issue with the word information alone.}
\end{classroomqa}

The Einstein--Podolsky--Rosen discomfort nevertheless survives. A singlet can give us maximal knowledge of the \emph{whole} while neither constituent has a pure local state. That reversal of the classical expectation---knowing the whole without assigning definite pure states to the parts---is not a signaling effect, but it is genuinely peculiar. To expose the peculiarity more sharply, turn quantum mechanics into a programming problem.

\section{A Classical Simulator for One Spin}

Draw a spin apparatus on a computer screen. The user can orient it, press a measurement button \(M\), and read \(+1\) or \(-1\). Inside the classical computer store two complex amplitudes,
\begin{equation}
\alpha_u,\qquad \alpha_d,
\qquad
|\alpha_u|^2+|\alpha_d|^2=1.
\label{eq:l8-one-spin-amplitudes}
\end{equation}
Between measurements, the computer integrates the appropriate Schr\"odinger equation. A simulated magnetic field can make the amplitudes precess exactly as a real spin state would.

\begin{figure}[htbp]
\centering
\includegraphics[width=0.88\textwidth]{lecture_08_figure_04.png}
\caption{The one-spin simulator: an orientable apparatus and measurement button connected to a computer storing two amplitudes.}
\lectureframe{00:27:47}
\end{figure}

When the user presses \(M\), a pseudorandom-number generator selects an outcome with Born probabilities
\begin{equation}
P(+1)=|\alpha_u|^2,
\qquad
P(-1)=|\alpha_d|^2.
\label{eq:l8-one-spin-born}
\end{equation}
The computer displays the selected sign and then performs the state update. After \(+1\),
\begin{equation}
(\alpha_u,\alpha_d)\longrightarrow(1,0),
\end{equation}
whereas after \(-1\),
\begin{equation}
(\alpha_u,\alpha_d)\longrightarrow(0,1).
\end{equation}
It then resumes unitary evolution from the post-measurement state. If the user rotates the apparatus before pressing \(M\), the computer changes basis and calculates the probabilities appropriate to the new axis.

Nothing specifically quantum is required of the hardware. A deterministic pseudorandom algorithm is enough to produce a sequence that looks random to the user. For a single spin, this classical imitation can be operationally indistinguishable from the laboratory experiment.

\begin{classroomqa}
\audiencequestion{What is the purpose of the measurement button? Does elapsed time matter before it is pressed?}
\lecturerresponse{The button marks the instant at which the simulator samples the Born distribution and updates the stored amplitudes. During the wait, the amplitudes evolve under the Hamiltonian. If the Hamiltonian is zero, nothing changes with time; with a magnetic field, the simulated spin precesses before the measurement.}
\end{classroomqa}

\begin{classroomqa}
\audiencequestion{Who decides when to press the button or which way to orient the apparatus?}
\lecturerresponse{In this model the experimentalist supplies those settings. The simulator continuously tracks the state, then uses the selected axis and measurement time when the button is pressed. The later discussion will return to the assumption that settings can be chosen independently of the prepared pair.}
\end{classroomqa}

\section{Two Spins: Where Local Simulation Fails}

Two spins in a product state remain easy. Give Alice one computer storing amplitudes \(\alpha_u,\alpha_d\), and Bob another storing \(\beta_u,\beta_d\). Each machine evolves and measures its own spin. No communication is needed because the joint coefficients factor:
\begin{equation}
\Psi(a,b)=\alpha_a\beta_b.
\end{equation}

An arbitrary two-spin state requires four coefficients,
\begin{equation}
\alpha_{uu},\qquad
\alpha_{ud},\qquad
\alpha_{du},\qquad
\alpha_{dd},
\qquad
\sum_{a,b}|\alpha_{ab}|^2=1.
\label{eq:l8-four-amplitudes}
\end{equation}
A single central computer can store these numbers and evolve them with the four-dimensional Schr\"odinger equation. When either observer measures, the central machine samples the correct marginal probability, discards the incompatible components, renormalizes the surviving components, and continues from the conditional state.

\begin{figure}[htbp]
\centering
\includegraphics[width=0.88\textwidth]{lecture_08_figure_05.png}
\caption{A central simulator stores the four amplitudes of an arbitrary two-spin state and remains connected to both measurement stations.}
\lectureframe{00:37:47}
\end{figure}

This construction simulates entanglement perfectly while the central processor remains connected to both sides. Trouble begins when one copies the four amplitudes into two machines, separates them, cuts the wires, and lets each machine use its own random generator. For the singlet,
\begin{equation}
\ket{\Psi^-}
=
\frac{1}{\sqrt2}
\left(
\ket{u}_A\ket{d}_B-\ket{d}_A\ket{u}_B
\right),
\label{eq:l8-singlet}
\end{equation}
equal-axis measurements must be opposite. Independent random choices can instead produce \(++\) or \(--\), so the correlation is lost even though both machines began with a copy of the joint coefficients.

\begin{classroomqa}
\audiencequestion{Could the two machines share one random-number source before they are separated?}
\lecturerresponse{Shared randomness can enforce a prescribed correlation for one fixed measurement axis. The deeper Bell problem asks for one local scheme that reproduces the quantum correlations for every independently chosen pair of axes. Preloaded shared variables are included among Bell-local models, and no such model reproduces the full quantum correlation function.}
\end{classroomqa}

\begin{classroomqa}
\audiencequestion{Are the four amplitudes independent?}
\lecturerresponse{They are four coordinates of one normalized joint state, not four independent local states. In an entangled state they cannot be written as \(\alpha_a\beta_b\). After a local result is selected, the surviving joint coefficients must be updated coherently; two disconnected copies do not automatically perform matching updates.}
\end{classroomqa}

\subsection{What the hidden wires mean}

A central random generator connected to both stations could coordinate the outputs. In an ordinary classical simulator those wires would also be communication channels. Bob could deliberately send a bit to Alice, and if the link were instantaneous the machine would violate locality.

The quantum system does not grant Bob that freedom. His allowed local transformations are unitary, and the reduced-density-matrix proof already showed that no such transformation changes Alice's marginal statistics. The simulator's wires are therefore a sign that ordinary classical bookkeeping has the wrong structure, not evidence that nature contains accessible superluminal cables.

\begin{classroomqa}
\audiencequestion{Do the wires required by the classical simulator represent physical nonlocality?}
\lecturerresponse{They represent the extra coordination a naive classical model needs. If they were ordinary usable wires, they would permit signaling and would not model the quantum experiment. The safer conclusion is that local classical simulation and quantum mechanics have different logical structures, not that hidden communication wires have been discovered.}
\end{classroomqa}

\begin{classroomqa}
\audiencequestion{If Alice and Bob each receive a random bit but the bits are perfectly correlated, what is strange about that?}
\lecturerresponse{Nothing is strange about one prearranged correlation; identical envelopes or shared random seeds can produce it. The nonclassical content appears when Alice and Bob independently choose among incompatible settings and the resulting correlations violate a Bell inequality while each local outcome remains random.}
\end{classroomqa}

The thought experiment also contains a premise that deserves daylight: the measurement settings are treated as independent choices. Alice and Bob decide when to press their buttons and how to orient their apparatuses without those choices being secretly fixed by the source variables. Once the observers and their decision mechanisms are included inside one closed physical system, the philosophical language of free choice becomes delicate. Operational Bell tests formulate the needed condition as measurement independence and arrange spacelike-separated, externally generated settings. The lecture leaves the metaphysics open while retaining the experimental premise.

\begin{classroomqa}
\audiencequestion{Does putting Alice and Bob inside the physical system create a loophole?}
\lecturerresponse{It exposes an assumption: the settings must not be correlated in advance with the hidden variables that prepare the pair. One may debate free will, but the testable statement is measurement independence. Abandoning it is logically possible, though it changes the explanatory burden rather than supplying an ordinary local simulation.}
\end{classroomqa}

\section{From Spins to a Particle on a Line}

Spins have carried the discussion from finite-dimensional vectors through entanglement and locality. Now enlarge the Hilbert space. Consider one particle moving on an infinite line. A position eigenstate is denoted
\begin{equation}
\ket{x},
\end{equation}
and means that the particle has definite position \(x\). These kets form a continuously labeled basis. A general state is a superposition of them, and its position-space wave function is
\begin{equation}
\psi(x)=\braket{x}{\Psi}.
\label{eq:l8-position-wavefunction}
\end{equation}

\begin{figure}[htbp]
\centering
\includegraphics[width=0.88\textwidth]{lecture_08_figure_06.png}
\caption{The position-space wave function as the projection of the abstract state onto the basis ket \(\ket{x}\).}
\lectureframe{00:56:51}
\end{figure}

Position is being used as a complete basis label. This does not say that a classical state can be specified by position alone, nor that the particle simultaneously has definite position and momentum. It says that every quantum state vector can be expanded in the position basis. Momentum will appear as an operator acting on those position-space coefficients.

The Born rule now gives a probability \emph{density},
\begin{equation}
\rho(x)=|\psi(x)|^2.
\label{eq:l8-probability-density}
\end{equation}
For a continuous variable, the probability of one exact point is ordinarily zero. The probability of finding the particle in an interval is
\begin{equation}
P(x_1<x<x_2)
=
\int_{x_1}^{x_2}|\psi(x)|^2\,dx,
\end{equation}
and normalization requires
\begin{equation}
\int_{-\infty}^{\infty}|\psi(x)|^2\,dx=1.
\label{eq:l8-continuous-normalization}
\end{equation}

\begin{classroomqa}
\audiencequestion{How can position alone be a complete basis when a classical particle needs both position and momentum?}
\lecturerresponse{Completeness here concerns the representation of a quantum vector. The coefficients \(\psi(x)\) for all \(x\) specify the state, and momentum information is encoded in their phases and spatial variation. Position and momentum are different bases, related later by a Fourier transform; they are not simultaneous classical labels.}
\end{classroomqa}

\begin{classroomqa}
\audiencequestion{Is \(|\psi(x)|^2\) the probability of finding the particle at the exact point \(x\)?}
\lecturerresponse{It is a density. Multiplying by a small interval \(dx\) gives the approximate probability in that interval, and integrating gives the exact probability for a finite region. A density may be nonzero even though the probability assigned to one isolated point is zero.}
\end{classroomqa}

\section{Continuous Orthonormality and the Delta Function}

For a discrete orthonormal basis,
\begin{equation}
\braket{a}{a'}=\delta_{aa'}.
\end{equation}
For a continuous basis the analogue is
\begin{equation}
\braket{x}{x'}=\delta(x-x'),
\label{eq:l8-delta-normalization}
\end{equation}
where the Dirac delta vanishes for \(x\neq x'\) and has unit area:
\begin{equation}
\int_{-\infty}^{\infty}\delta(x-x')\,dx'=1.
\end{equation}
One may picture it as a narrow spike of width \(\epsilon\), height \(1/\epsilon\), and fixed area, followed by the limit \(\epsilon\to0\).

\begin{figure}[htbp]
\centering
\includegraphics[width=0.88\textwidth]{lecture_08_figure_07.png}
\caption{The board's narrow-spike representation of the Dirac delta: width \(\epsilon\), height \(1/\epsilon\), and unit area.}
\lectureframe{01:01:57}
\end{figure}

The notation is distributional shorthand. A position eigenket is not a normalizable vector in the ordinary square-integrable Hilbert space; it is a generalized vector whose delta normalization is meaningful under an integral. The central working rule is
\begin{equation}
\int_{-\infty}^{\infty}
\delta(x-x')f(x')\,dx'
=f(x).
\label{eq:l8-delta-sampling}
\end{equation}
The delta selects the value for which its argument vanishes.

Now write the continuous expansion
\begin{equation}
\ket{\Psi}
=
\int_{-\infty}^{\infty}
\psi(x')\ket{x'}\,dx'.
\label{eq:l8-continuous-expansion}
\end{equation}
Acting with \(\bra{x}\) gives
\begin{align}
\braket{x}{\Psi}
&=
\int dx'\,\psi(x')\braket{x}{x'} \notag\\
&=
\int dx'\,\psi(x')\delta(x-x') \notag\\
&=
\psi(x).
\label{eq:l8-expansion-consistency}
\end{align}
Thus the coefficient in the expansion and the projection \(\braket{x}{\Psi}\) are the same function, exactly as in a discrete basis.

\begin{figure}[htbp]
\centering
\includegraphics[width=0.88\textwidth]{lecture_08_figure_08.png}
\caption{The continuous state expansion and its consistency check using \(\braket{x}{x'}=\delta(x-x')\).}
\lectureframe{01:07:00}
\end{figure}

\begin{classroomqa}
\audiencequestion{Does \(\psi(x')\) play the same role as the amplitudes in a finite-dimensional vector?}
\lecturerresponse{Yes. The basis label \(a'\) becomes the continuous label \(x'\), the sum becomes an integral, and the coefficient \(\psi(a')\) becomes the function \(\psi(x')\). The underlying vector-space statement is unchanged.}
\end{classroomqa}

\begin{classroomqa}
\audiencequestion{Is the Dirac delta really a function with an infinite value at one point?}
\lecturerresponse{That picture is useful but not literal. Rigorously, the delta is a distribution defined by its action under an integral, Equation~\eqref{eq:l8-delta-sampling}. The narrow unit-area spike is a sequence of ordinary functions that approaches the same sampling rule.}
\end{classroomqa}

\section{Functions as Vectors}

To see the inner product, first recall the discrete calculation:
\begin{align}
\braket{\psi}{\phi}
&=
\sum_{a',a}
\psi^*(a')\braket{a'}{a}\phi(a) \notag\\
&=
\sum_a\psi^*(a)\phi(a).
\label{eq:l8-discrete-inner-product}
\end{align}
The complex-conjugated coefficients belong to the bra, and orthonormality collapses the double sum.

\begin{figure}[htbp]
\centering
\includegraphics[width=0.88\textwidth]{lecture_08_figure_09.png}
\caption{The discrete inner product expanded in a basis before sums are replaced by integrals.}
\lectureframe{01:11:11}
\end{figure}

Replace the basis labels by positions and the sum by an integral:
\begin{equation}
\braket{\psi}{\phi}
=
\int_{-\infty}^{\infty}\psi^*(x)\phi(x)\,dx.
\label{eq:l8-continuous-inner-product}
\end{equation}
Interchanging the vectors complex conjugates the answer,
\begin{equation}
\braket{\phi}{\psi}
=
\braket{\psi}{\phi}^*,
\end{equation}
and the norm is
\begin{equation}
\|\psi\|^2=\braket{\psi}{\psi}
=\int |\psi(x)|^2\,dx.
\end{equation}

Functions are therefore vectors: they can be added, multiplied by complex scalars, normalized, and paired by an inner product. Physical wave functions on the infinite line are square-integrable, so functions that grow without bound at infinity are excluded. The generalized basis functions \(\delta(x-x_0)\) sit just outside that Hilbert space but remain indispensable as distributional basis vectors.

\begin{classroomqa}
\audiencequestion{Should the unmarked integration limits be \(-\infty\) and \(+\infty\)?}
\lecturerresponse{Yes, for the infinite line. Those limits also explain why normalizable wave functions must fall off sufficiently fast: otherwise the total probability integral would not exist.}
\end{classroomqa}

\begin{classroomqa}
\audiencequestion{A vertical bar was missing in one of the board expressions. Does that change the meaning?}
\lecturerresponse{It was only a notation slip. The basis overlap is \(\braket{a'}{a}\), with a bra on the left and a ket on the right. Restoring the bar makes clear that the expression is an inner product rather than an ambiguous row or column.}
\end{classroomqa}

\begin{classroomqa}
\audiencequestion{Is almost everything entangled with an environment most of the time?}
\lecturerresponse{In realistic macroscopic conditions, yes. Systems are continually struck by photons, air molecules, and other degrees of freedom, spreading phase information into the environment. That environmental entanglement is central to decoherence and to why interference between macroscopic alternatives is hard to observe.}
\end{classroomqa}

\begin{classroomqa}
\audiencequestion{Did the mathematical theory of distributions change physics, or merely make the delta function respectable?}
\lecturerresponse{Its immediate service was to make rigorous manipulations that physicists were already using. It supplied a precise language for generalized functions and reduced arguments with mathematicians, although it did not by itself resolve every distributional difficulty in quantum field theory.}
\end{classroomqa}

\section{Linear Operators on Wave Functions}

A linear operator maps one function to another while respecting superposition:
\begin{equation}
L(c_1\psi+c_2\phi)
=
c_1L\psi+c_2L\phi.
\end{equation}
An observable must be represented by a Hermitian operator. In the present notation the condition is
\begin{equation}
\braket{\phi}{L\psi}
=
\braket{L\phi}{\psi},
\label{eq:l8-hermitian-condition}
\end{equation}
or equivalently
\begin{equation}
\braket{\phi}{L\psi}
=
\braket{\psi}{L\phi}^*.
\end{equation}
For unbounded operators, this statement includes a domain and boundary conditions; the integrations below assume wave functions for which the boundary terms vanish.

\subsection{Position by multiplication}

The simplest operator multiplies a wave function by the coordinate:
\begin{equation}
(x\psi)(x)=x\psi(x).
\label{eq:l8-position-operator}
\end{equation}
It is linear, and
\begin{align}
\braket{\phi}{x\psi}
&=
\int dx\,\phi^*(x)x\psi(x), \notag\\
\braket{x\phi}{\psi}
&=
\int dx\,[x\phi(x)]^*\psi(x) \notag\\
&=
\int dx\,\phi^*(x)x\psi(x),
\end{align}
because \(x\) is real. Multiplication by \(x\) is therefore Hermitian and represents the position observable.

\begin{figure}[htbp]
\centering
\includegraphics[width=0.88\textwidth]{lecture_08_figure_10.png}
\caption{The Hermiticity check for the position operator, viewed as multiplication of a wave function by the real variable \(x\).}
\lectureframe{01:26:47}
\end{figure}

\begin{classroomqa}
\audiencequestion{In \(x\psi(x)\), is \(x\) an operator or merely a real number?}
\lecturerresponse{At each point it is the real coordinate value, but the rule that maps the whole function \(\psi(x)\) to the new function \(x\psi(x)\) is the operator. Its reality is exactly what makes the two inner products in the Hermiticity test agree.}
\end{classroomqa}

\subsection{Differentiation is anti-Hermitian}

Try the derivative
\begin{equation}
D=\frac{d}{dx}.
\end{equation}
It is linear, but integration by parts gives
\begin{align}
\braket{\phi}{D\psi}
&=
\int dx\,\phi^*(x)\frac{d\psi}{dx} \notag\\
&=
\left[\phi^*(x)\psi(x)\right]_{-\infty}^{\infty}
-
\int dx\,\frac{d\phi^*}{dx}\psi(x) \notag\\
&=
-\braket{D\phi}{\psi},
\label{eq:l8-derivative-adjoint}
\end{align}
where the boundary term vanishes. Thus
\begin{equation}
D^\dagger=-D.
\end{equation}
The derivative is anti-Hermitian rather than Hermitian.

\begin{figure}[htbp]
\centering
\includegraphics[width=0.88\textwidth]{lecture_08_figure_11.png}
\caption{Integration by parts moves the derivative between the two wave functions and introduces the decisive minus sign.}
\lectureframe{01:32:30}
\end{figure}

Multiply by \(-i\):
\begin{equation}
p=-i\frac{d}{dx}.
\label{eq:l8-momentum-natural-units}
\end{equation}
Complex conjugation changes \(-i\) to \(+i\), while integration by parts contributes another minus sign. The signs cancel, so \(p^\dagger=p\). Restoring units,
\begin{equation}
p=-i\hbar\frac{d}{dx}.
\label{eq:l8-momentum-operator}
\end{equation}

The expectation value makes its reality explicit:
\begin{align}
\expect{p}_\psi
&=
-i\int dx\,\psi^*(x)\frac{d\psi}{dx} \notag\\
&=
i\int dx\,\frac{d\psi^*}{dx}\psi(x) \notag\\
&=
\expect{p}_\psi^*.
\label{eq:l8-momentum-expectation-real}
\end{align}
Any number equal to its own complex conjugate is real, as an observable's expectation value must be.

\begin{figure}[htbp]
\centering
\includegraphics[width=0.88\textwidth]{lecture_08_figure_12.png}
\caption{The momentum expectation value before and after integration by parts; the two forms are complex conjugates and equal.}
\lectureframe{01:39:44}
\end{figure}

\begin{classroomqa}
\audiencequestion{Why can an operator containing \(i\) represent a real observable?}
\lecturerresponse{Hermiticity concerns the full operator under the inner product, not whether every symbol in its coordinate formula is real. The factor \(i\) compensates for the minus sign acquired when the derivative is moved by integration by parts, making all expectation values real.}
\end{classroomqa}

\section{The Canonical Commutator}

One more calculation reveals why the operator in Equation~\eqref{eq:l8-momentum-operator} deserves attention. In units with \(\hbar=1\), act first with \(x\) and then with \(p\):
\begin{align}
(px\psi)(x)
&=
-i\frac{d}{dx}\bigl(x\psi(x)\bigr) \notag\\
&=
-i\psi(x)-ix\frac{d\psi}{dx}.
\label{eq:l8-px}
\end{align}
Reverse the order:
\begin{equation}
(xp\psi)(x)
=
-ix\frac{d\psi}{dx}.
\label{eq:l8-xp}
\end{equation}
Subtracting Equation~\eqref{eq:l8-px} from Equation~\eqref{eq:l8-xp},
\begin{equation}
\bigl([x,p]\psi\bigr)(x)
=
\bigl(xp-px\bigr)\psi(x)
=
i\psi(x).
\end{equation}
Because this holds for every wave function in the operator domain,
\begin{equation}
[x,p]=i,
\end{equation}
or, with units restored,
\begin{equation}
[x,p]=i\hbar.
\label{eq:l8-canonical-commutator}
\end{equation}

\begin{figure}[htbp]
\centering
\includegraphics[width=0.88\textwidth]{lecture_08_figure_13.png}
\caption{The completed operator argument: if \(([x,p]-i)\psi=0\) for every wave function, then \([x,p]=i\).}
\lectureframe{01:46:24}
\end{figure}

The classical counterpart is
\begin{equation}
\{x,p\}_{\mathrm{PB}}=1.
\end{equation}
Dirac's historical reasoning ran in the opposite direction from the present derivation. He began with the algebraic resemblance between antisymmetric Poisson brackets and commutators, postulated the correspondence
\begin{equation}
[A,B]\longleftrightarrow i\hbar\{A,B\}_{\mathrm{PB}},
\end{equation}
and inferred that the momentum operator must satisfy the canonical commutator with position. Here we discovered \(-i\,d/dx\) by repairing the derivative's failure of Hermiticity and only afterward found the same algebra.

\begin{classroomqa}
\audiencequestion{Does the commutator alone prove that \(-i\hbar\,d/dx\) is physical momentum?}
\lecturerresponse{It is strong evidence, not the whole physical identification. We must still show that the operator generates spatial translations, that its expectation value is conserved when translation symmetry is present, and that a wave packet moves with group velocity \(\expect{p}/m\) in the appropriate nonrelativistic system.}
\end{classroomqa}

\section{What Comes Next}

The move from spins to particles has preserved the architecture while changing its scale. A continuously labeled basis replaces a finite list of basis vectors. Sums become integrals, Kronecker deltas become Dirac deltas, and column vectors become square-integrable functions. Position acts by multiplication, while momentum acts by differentiation.

The next tasks are physical rather than merely notational. Momentum probabilities require the momentum basis and Fourier transformation. Time evolution requires the Schr\"odinger equation for \(\psi(x,t)\). Wave packets will then connect the abstract operator \(p=-i\hbar\,d/dx\) to motion: their group velocity will be related to momentum divided by mass, and translation symmetry will explain momentum conservation. The commutator \([x,p]=i\hbar\) is therefore not an isolated trick. It is the hinge between the vector-space mechanics already developed and the quantum dynamics of particles.
